\part{Galders: Poetic Charms, Spells, and Curses}

\bookStart{Introduction to Galders}

The Germanic word “galder” (PGmc. \emph{*galdrą}) is a verbal noun derived from the old verb “gale” (PGmc. \emph{*galaną}), meaning ‘to sing, to chant’.  A “galder” is therefore what is chanted or sung, more specifically a magical chant, whether for the purpose of healing or cursing.

In the present section the reader will find gathered a number of such \inx[C]{galder}[galders] attested across the Old Germanic languages.  Not all of these texts may be true galders, but they share a magical significance and many formulae and stylistic features occur throughout the corpus.

This section will only include written charms with clear pre-Christian or otherwise traditional elements like the invocation of the Heathen Gods or use of poetic formulae.  Some texts are syncretic and mix traditional elements with invocations of Jesus Christ and Saint Mary, but thoroughly Christian prayers, even if metrical, are excluded from this section and will instead be found under “Poetry on Christian Subjects”.


\chapter{Continental Germanic galders}

\input{books/galders/Merseburg.tex}
\newpage
\input{books/galders/Against Worms.tex}


\chapter{Old English galders}

\bookStart{Acre-Boot}[Æcer-bót]
\setBookCode{Acerbot}

\begin{flushright}%
\textbf{Dating:} ca. 1000

\textbf{Meter:} \Fornyrdislag%
\end{flushright}%

\section{Introduction}

The \textbf{Acre-Boot} (i.e.~‘Field-Cure’, abbrev. \Acerbot) is an unusually long text mixing Catholic, Germanic, and Celtic elements.

\Acerbot\ is preserved in Cotton MS Caligula A VII, foll.~176r--178r.

\sectionline

\section{Text}

\bpg\bpa Hér ys seo bót, hú þú meaht þíne æceras bétan gif hí nellaþ wel wexan oþþe þǽr hwilc un·ge·défe þing on ge·dón bið on dry oððe on lyb-lâce. Ge·nim þǫnne ǫn niht, ær hyt dagige, feower tyrf on feower healfa þæs landes, and ge·mearca hú hý ǽr stódon. Nim þonne ele and hunig and beorman, and ælces feos meolc þe on þæm lande sý, and ælces treow-cynnes dǽl þe on þǽm lande sý ge·wexen, bútan heardan béaman, and ælcre nam cúþre wyrte dǽl, bútan glappan ânon, and dó þonne hâlig-wæter ðǽr on, and drype þǫnne þriwa on þǫne staðol þára turfa, and cweþe ðonne ðás word:
„Crescite, wexe, et multiplicamini, and ge·mænig fealda, et replete, and ge·fylle, terre, þás eorðan.
In nomine patris et filii et spiritus sancti sit benedicti.“
And Pater Noster swá oft swá þæt óðer.  And bere siþþan ðá turf to circean, and mæsse-preost á·singe feower mæssan ofer þan turfon, and wende man þæt gréne to ðan weofode, and siþþan ge·bringe man þá turf þǽr hí ǽr wǽron ǽr sunnan setl-gange.
And hæbbe him gæ·worht of cwic-béame feower Cristes mælo and a·wríte on ælcon ende:
Matheus and Marcus, Lucas and Iohannes.
Lege þæt Cristes mæl on þone pyt neoþe-weardne, cweðe ðǫnne:
Crux Matheus, crux Marcus, crux Lucas, crux Sanctus Iohannes.\epa

\bpb TODO.\epb\epg


\bpg\bpa Nim ðǫnne þá turf and sete ðær ufon on and cweþe ðǫnne nigon síþon þás word:
Crescite, and swa oft Pater Noster, and wende þé þǫnne éast-weard, and on·lút nigon síðon éad-mód⸗líce, and cweð þǫnne þás word:\epa

\bpb TODO.\epb\epg


\bvg\bva%
„\alst{Éa}st-weard ic stande, \hld\ \alst{â}rena ic mé bidde, &
bidde ic þǫne \alst{m}ǽran domine, \hld\ bidde ðǫne \alst{m}iclan drihten, &
bidde ic ðǫne \alst{h}áligan \hld\ \alst{h}eofon-rices weard, &
\alst{eo}rðan ic bidde \hld\ and \alst{u}p-heofon &
and ðá \alst{s}óþan \hld\ \alst{s}ancta Marian &
and \alst{h}eofones meaht \hld\ and \alst{h}éah-reced, &
þæt ic móte þis \alst{g}ealdor \hld\ mid \alst{g}ife drihtnes &
\alst{t}óðum on·\alst{t}ýnan \hld\ þurh \alst{t}rumne ge·þanc, &
á·\alst{w}ęccan þas \alst{w}æstmas \hld\ ús to \alst{w}oruld-nytte, &
ge·\alst{f}ylle þás \alst{f}oldan \hld\ mid \alst{f}æste ge·léafan &
\alst{w}litigigan þas \alst{w}ancg-turf, \hld\ swá sé \alst{w}itega cwæð &
þæt sé hæfde \alst{á}re ǫn \alst{eo}rþ-ríce, \hld\ sé þe \alst{æ}lmyssan &
\alst{d}ǽlde \alst{d}óm-líce \hld\ \alst{d}rihtnes þances.“\eva

\bvb TODO.\evb\evg


\bpg\bpa Wende þe þǫnne III sun-ganges, a·strece þǫnne ǫn and-lang and a·rim þǽr letanias and cweð þonne: „Sanctus, sanctus, sanctus oþ ende.“\epa

\bpb TODO.\epb\epg


\bpg\bpa Sing þǫnne Benedicite a·þenedon earmon and Magnificat and Pater Noster III, and be·béod hit Criste and sancta Marian and þǽre hâlgan róde tó lofe and to weorþinga and þâm âre þe þæt land âge and eallon þâm þe him under·ðeodde synt.\epa

\bpb TODO.\epb\epg


\bpg\bpa Ðǫnne þæt eall síe ge·dón, þǫnne nime man un·cúþ sǽd æt ælmes-mannum and selle him twá swylc, swylce man æt him nime, and ge·gaderie ealle his sulh-ge·teogo to·gædere; borige þǫnne on þam béame stor and finol and ge·hâlgode sápan and ge·hâlgod sealt.
Nim þǫnne þæt sǽd, sete on þæs sules bodig, cweð þǫnne:\epa

\bpb TODO.\epb\epg


\bvg\bva%
„\alst{E}rce, \alst{E}rce, \alst{E}rce, \hld\ \alst{eo}rþan módor, &
ge·\alst{u}nne þé se \alst{a}l-walda, \hld\ \alst{é}ce drihten, &
æcera \alst{w}exendra \hld\ and \alst{w}rídendra, &
\alst{ea}cniendra \hld\ and \alst{e}lniendra; &
\alst{sc}eafta héhre \hld\ \alst{sc}íre wæstma, &
and þǽre \alst{b}râdan \hld\ \alst{b}ere-wæstma, &
and þǽre \alst{h}wítan \hld\ \alst{h}wæte-wæstma, &
and \alst{ea}lra \hld\ \alst{eo}rþan wæstma.\eva

\bvb TODO.\evb\evg


\bvg\bva%
Ge·\alst{u}nne him \hld\ \alst{é}ce drihten &
and his \alst{h}âlige, \hld\ þe ǫn \alst{h}eofonum synt, &
þæt hys yrþ sí ge·\alst{f}riþod \hld\ wið ealra \alst{f}eonda ge·hwæne, &
and heo sí ge·\alst{b}orgen \hld\ wið ealra \alst{b}ealwa ge·hwylc, &
þára \alst{l}yb-lâca \hld\ geond \alst{l}and sáwen.\eva

\bvb TODO.\evb\evg


\bvg\bva%
Nú ic bidde ðǫne \alst{w}aldend, \hld\ sé ðe ðás \alst{w}or-uld ge·sceóp, &
þæt ne sý nán tó þæs \alst{c}widol wíf \hld\ ne tó þæs \alst{c}ræftig man &
þæt á·\alst{w}endan ne mæge \hld\ \alst{w}ord þus ge·cwedene.“\eva

\bvb TODO.\evb\evg



\bpg\bpa Þonne man þá sulh forð drife þá forman furh on·scéote, cweð þǫnne:\epa

\bpb TODO.\epb\epg


\bvg\bva%
„Hál wes þú, \alst{F}olde, \hld\ \alst{f}ir\emph{h}a módor! &
Beo þú \alst{g}rówende \hld\ on \alst{G}odes fæþme, &
\alst{f}ódre ge·\alst{f}ylled \hld\ \alst{f}ir\emph{h}um tó nytte.“\eva

\bvb “Hale be thou, Fold, mother of men! \\
Be thou growing in God’s bosom, \\
filled with fodder as a boon to men.”\evb\evg


\bpg\bpa Nim þonne ælces cynnes melo and a·bacæ man inn-werdre handa brádnæ hláf and ge·cned hine mid meolce and mid hálig-wætere and lecge under þa forman furh. Cweþe þonne:\epa

\bpb TODO.\epb\epg


\bvg\bva%
„\alst{F}ul æcer \alst{f}ódres \hld\ \alst{f}ir\emph{h}a cinne, &
\alst{b}eorht-\alst{b}lowende, \hld\ þú ge·\alst{b}létsod weorþ &
þæs \alst{h}âligan nǫman \hld\ þe ðás \alst{h}eofon ge·sceop &
and þás \alst{eo}rþan \hld\ þe wé \alst{ǫ}n lifiaþ; &
sé \alst{g}od, sé þás \alst{g}rundas ge·worhte, \hld\ ge·unne ús \alst{g}rowende gife, &
þæt ús \alst{c}orna ge·hwylc \hld\ \alst{c}ume tó nytte.“\eva

\bvb “Acre, full of fodder for the race of men, \\
brightly blooming; thou wast blessed \\
by the holy name which created the heaven \\
and this earth upon which we live. \\
The God who made the lands, grant us growing grace \\
that each kind of grain may be for us a boon.”\evb\evg


\bpg\bpa Cweð þonne: „III Crescite in nomine patris, sit benedicti.“  Amen and Pater Noster þríwa.\epa

\bpb TODO.\epb\epg

\newpage
\input{books/galders/Against Swarm.tex}
\newpage
\input{books/galders/Against Dwarf.tex}
\newpage
\input{books/galders/Against Fearstitch.tex}
\newpage
\bookStart{The Galder of Nine Worts}[Nigon wyrta galdor]
\setBookCode{NineHerbs}

\begin{flushright}%
\textbf{Dating:} ?

\textbf{Meter:} \Fornyrdislag%para
\end{flushright}%

\section{Introduction}

TODO: introduction

Numerous lines begin with small crosses which appear to serve as section markers.  The present division into stanzas is largely based on these.  TODO: note which.

Preserved in \Lacnunga.

\sectionline

\section{Text}

\bvg\bva%
Ge·\alst{m}yne ðú \alst{m}ug-wyrt \hld\ hwæt þú á·\alst{m}eldodest &
hwæt þu \alst{r}enadest \hld\ æt \alst{R}egen-melde?\eva

\bvb Rememberest thou, Mugwort, what thou didst declare, \\
what thou didst arrange at Reinmeld?\evb\evg


\bvg\bva \alst{U}na þú hâttest \hld\ \alst{y}ldost wyrta &
þú miht \edtrans{wið \emph{\alst{þ}rim} \hld\ and wið \emph{\alst{þ}rí·tigum}}{against three and against thirty}{\Bfootnote{I.e. ‘against a great number of foes’.  ‘Three and thirty’ is formulaic and appears in many later English and Scottish folk-songs, viz. Child 4EFG, 18B, 20C, 30, 53BCDEIKM, 63EFH, 73I, 97AC, 100AG, 110BGH, 156G, 185A, 187A, 187C, 190A, 192A, 193B, 203C, 211A, 217GHLN, 244A, 268A, 269C, 281ABC.  Things described include horses, heads of cattle, warriors, days, years, winters.}} &
þú miht wiþ \alst{a}ttre \hld\ and wið \alst{ǫ}n·flyge &
þú miht wiþ þâm \alst{l}âþan \hld\ ðe geond \alst{l}ǫnd færð.\eva

\bvb Un thou art called, the oldest of worts; \\
thou availest against three and against thirty; \\
thou availest against the venom and against the onflier; \\
thou availest against the loathsome one that journeys through the lands.\evb\evg


\bvg\bva Ond þú \alst{w}eg·bráde \hld\ \alst{w}yrta módor &%NOTE: introduced with a + sign
\alst{éa}stan \alst{o}pene \hld\ \alst{i}nnan mihtigu &
ofer ðy \alst{c}ræte \alst{c}urron \hld\ ofer ðy \alst{c}wéne reodon &
\ind ofer ðy \alst{b}rýde \alst{b}rýodedon &
\ind ofer ðy \alst{f}earras \alst{f}nærdon.\eva

\bvb And thou, Waybroad, mother of worts, \\
open from the east, mighty within. \\
Over thee TODO. \\
\ind TODO \\
\ind TODO\evb\evg


\bvg\bva Eallum þu þon wið·\alst{st}óde \hld\ and wið·\alst{st}unedest &
swá ðú wið·stonde \alst{a}ttre \hld\ and \alst{ǫ}n·flyge &
and þǽm \alst{l}âðan \hld\ þe geond \alst{l}ǫnd fereð.\eva

\bvb Them all didst thou then withstand, and didst stop; \\
so mayst thou withstand the venom and the onflier, \\
and the loathsome one that journeys through the lands.\evb\evg


\bvg\bva \alst{St}une hætte þéos wyrt, \hld\ héo on \alst{st}âne ge·weox &
\alst{st}ǫnd héo wið attre, \hld\ \alst{st}unað héo wærce &
\alst{St}iðe héo hatte, \hld\ wið·\alst{st}unað héo attre &
\alst{w}receð héo \alst{w}râðan, \hld\ \alst{w}eorpeð út attor.\eva

\bvb Stun is this wort called, she grew on stone; \\
she withstands venom, she stops aches. \\
Stithe is she called, she stops the venom; \\
she drives away the wroth one, casts out the venom.\evb\evg


\bvg\bva Þis is séo \alst{w}yrt \hld\ séo wiþ \alst{w}yrm ge·feaht &
þéos mæg wið \alst{a}ttre, \hld\ héo mæg wið \alst{ǫ}n·flyge; &
héo mæg wið ðâm \alst{l}âþan \hld\ ðe geond \alst{l}ǫnd fereþ.\eva

\bvb This is the wort that fought against the Wyrm; \\
this one avails against the venom, she avails against the onflier; \\
she avails against the loathsome one that journeys through the lands.\evb\evg


\bvg\bva Fleoh þú nú attor-\alst{l}âðe, \hld\ séo \alst{l}ǽsse ðá mâran &
séo mâre þá lǽssan, \hld\ oððæt him beigra bót sý!\eva

\bvb TODO \\
TODO.\evb\evg


\bvg\bva Ge·\alst{m}yne þú, \alst{m}ægðe, \hld\ hwæt þú á·\alst{m}eldodest &
hwæt ðú ge·\alst{æ}ndadest \hld\ æt \alst{A}lor-forda &
þæt nǽfre for ge·\alst{f}loge \hld\ \alst{f}eorh ne ge·sealde &
syþðan him \alst{m}ǫn \alst{m}ægðan \hld\ tú \alst{m}ęte ge·gyrede.\eva

\bvb TODO \\
TODO \\
TODO \\
TODO\evb\evg


\bvg\bva Þis is séo \alst{w}yrt \hld\ ðe \alst{w}er-gulu hatte &
ðás on·\alst{s}ænde \alst{s}eolh \hld\ ofer \alst{s}ǽs hrygc &
\alst{o}ndan \alst{a}ttres \hld\ \alst{ȯ}þres tó bóte.\eva

\bvb TODO \\
TODO \\
TODO\evb\evg


\bvg\bva Ðás \emph{\alst{n}ygon} magon \hld\ wið \alst{n}ygon attrum.\eva

\bvb These nine avail against nine venoms.\evb\evg


\bvg\bva Wyrm cóm \alst{s}nícan, \hld\ to·\alst{s}lât hé \emph{m}an &
ðá ge·nam \alst{W}óden \hld\ \emph{\alst{n}ygon} \alst{w}uldor-tânas &
slóh ðá þá \alst{n}ǽddran \hld\ þæt héo on \emph{\alst{n}ygon} tó·fléah &
Þǽr ge·\alst{æ}ndade \alst{æ}ppel \hld\ and \alst{a}ttor &
þæt héo nǽfre ne wolde \hld\ ǫn hús búgan.\eva

\bvb A \inx[C]{Wyrm} came crawling; he tore apart a man. \\
Then took Weden nine glory-twigs, \\
slew then that adder, that it sprung into nine [parts]. \\
There ended apple and venom, \\
that she would never wish to enter a house.\evb\evg


\bvg\bva \alst{F}ille and \alst{f}inule, \hld\ \alst{f}ela-mihtigu twá &
þá \alst{w}yrte ge·sceop \hld\ \alst{w}ítig drihten &
\alst{h}âlig on \alst{h}eofonum, \hld\ þá hé \alst{h}ǫngode; &
\alst{s}ętte and \alst{s}ęnde \hld\ on \emph{\alst{s}eofon} worulde &
\alst{ea}rmum and \alst{éa}digum \hld\ \alst{ea}llum tó bóte.\eva

\bvb Fill and Fennel, the many-mighty two; \\
those worts the wise lord shaped, \\
holy in heaven when he hung. \\
He set and sent them into seven worlds, \\
for wretched men and for wealthy, for all men as a cure.\evb\evg


\bvg\bva \alst{St}ǫnd héo wið wærce, \hld\ \alst{st}unað héo wið attre &
séo mæg wið \emph{\alst{þ}rim} \hld\ \emph{and} wið \emph{\alst{þ}rí·tigum} &
wið \emph{\alst{f}éondes} hǫnd \hld\ and wið \alst{f}ǽr-bregde &
wið \alst{m}alscrunge \hld\ \alst{m}ânra wihta.\eva

\bvb She stands against ache, she stands against venom;
she avails against three and against thirty;
against the TODO of evil wights.\evb\evg


\bvg\bva Nu magon þás \emph{nygon} \alst{w}yrta \hld\ wið nygon \alst{w}uldor-ge·flogenum &
wið \emph{nygon} \alst{a}ttrum \hld\ and wið nygon \alst{o}n·flygnum &
wið ðý \alst{r}éadan attre, \hld\ wið ðý \alst{r}unlan attre &
wið ðý \alst{h}witan attre, \hld\ wið ðý \emph{\alst{h}æwe}nan attre &
wið ðý \alst{g}eolwan attre, \hld\ wið ðý \alst{g}rénan attre &
wið ðý \alst{w}onnan attre, \hld\ wið ðý \alst{w}edenan attre &
wið ðý \alst{b}rúnan attre, \hld\ wið ðý \alst{b}asewan attre &
wið \alst{w}yrm-ge·blæd, \hld\ wið \alst{w}æter-ge·blæd &
wið \alst{þ}orn-ge·blæd, \hld\ wið \alst{þ}ystel-ge·blæd &
wið \alst{ý}s-ge·blæd, \hld\ wið \alst{a}ttor-ge·blæd\eva

\bvb Now these nine worts avail against glory-onfliers: \\
against nine venoms and against nine onfliers; \\
against the red venom; against the TODO venom; \\
against the white venom; against the TODO venom; \\
against the yellow venom; against the green venom; \\
against the TODO venom; against the TODO venom; \\
against the brown venom; against the TODO venom; \\
against worm-TODO; against water-TODO; \\
against thorn-TODO; against thistle-TODO; \\
against ice-TODO; against venom-TODO.\evb\evg


\bvg\bva Gif ænig \alst{a}ttor cume \hld\ \alst{éa}stan fleógan &
oððe ǽnig norðan cume &
oððe ǽnig \alst{w}estan \hld\ ofer \alst{w}er-þeóde\eva

\bvb If any venom should come flying from the east; \\
or any come from the north; \\
or any from the west, over mankind.\evb\evg


\bvg\bva Críst stód ofer \alst{á}dle \hld\ \alst{ǽ}ngan cundes. &
Ic \alst{â}na wât \hld\ \alst{éa} rinnende &
þǽr þá \alst{n}ygon \alst{n}ǽdran \hld\ \alst{n}éan be·healdað.\eva

\bvb Christ stood over TODO; \\
I know one river running, \\
where the nine adders TODO.\evb\evg


\bvg\bva Motan ealle \alst{w}éoda \hld\ nu \alst{w}yrtum á·springan &
\alst{s}ǽs tó·\alst{s}lúpan, \hld\ eal \alst{s}ealt wæter &
ðonne ic \alst{þ}is attor \hld\ of \alst{þ}é ge·bláwe.\eva

\bvb TODO\evb\evg


\bpg\bpa Mucgwyrt, weg-brade þe eastan open sy, lombes-cyrse, attor-laðan, mageðan, netelan, wudu-sur-æppel, fille and finul, ealde sapan. Ge·wyrc ða wyrta to duste, mængc wiþ þa sapan and wiþ þæs æpples gor. Wyrc slypan of wætere and of axsan, ge·nim finol, wyl on þære slyppan and beþe mid æggemongc, þonne he þa sealfe on do, ge ær ge æfter. Sing þæt galdor on æcre þara wyrta, :III: ær he hy wyrce and on þone æppel eal-swa; ond singe þon męn in þǫne mu̇ð and in þá éaran búta and on ðá wunde þæt ilce gealdor, ær he þá sealfe on dó.\epa

\bpb TODO.\epb\epg

\sectionline



\chapter{Old Norse galders}

\bookStart{Ribe galder stick (DR EM85;493)}
\setBookCode{RibeStick}

\begin{flushright}%
\textbf{Dating:} Mediæval.%TODO

\textbf{Meter:} \Fornyrdislag, \Galdralag%para
\end{flushright}%

\subsection{Introduction}

A wooden stick from the Danish city of Ribe.  The galder is syncretic and contains numerous pre-Christian elements in a Christian(ised) context.

The inscription may be conveniently divided into four parts.  Part one (ll.~1--4) contains an introductory prayer where the healer asks for the aid of natural forces (Earth, Up-heaven and the Sun) and Christian divinitities (God and Saint Mary) so that the healing may be successful.  Part two (ll.~5--8) ritually exorcises any sickness which may have entered any part of the body.  Part three (ll.~9--14) apparently warns the addressee that they will be haunted by “nine needs” (an old Heathen formula; see note) until they say the charm.  Part four (ll. 15, which is probably prose) gives the personal name “Bonde”, perhaps the addressee, and concludes with an “Amen”.

\subsection{Transliteration}

Note that the part of § C following ¶ is deliberately scraped off.

§ A \textbf{+ io\textasciicircum rþ ÷ biþ a\textasciicircum k ÷ ua\textasciicircum rþæ ÷ o\textasciicircum k ÷ uphimæn ÷ so\textasciicircum l ÷ o\textasciicircum k ÷ sa\textasciicircum nt\textasciicircum æ maria ÷ o\textasciicircum k ÷ salfæn ÷ gud| |drotæn ÷ þæt han ÷ læ mik ÷ læknæs÷ha\textasciicircum nd ÷ o\textasciicircum k lif÷tuggæ ÷ at\textasciicircum  \textasciicircum liuæ} \\
§ B \textbf{uiuindnæ ÷ þær ÷ botæ ÷ þa\textasciicircum rf ÷ or ÷ ba\textasciicircum k ÷ o\textasciicircum k or brʀst ÷ or lækæ ÷ o\textasciicircum k or lim ÷ or øuæn ÷ o\textasciicircum k or øræn ÷ or ÷ a\textasciicircum llæ þe ÷ þær ÷ ilt ÷ kan i at} \\
§ C \textbf{kumæ ÷ suart ÷ hetær ÷ sten ÷ ha\textasciicircum n ÷ stær ÷ i ÷ hafæ ÷ utæ ÷ þær ÷ ligær ÷ a ÷ þe ÷ ni ÷ no\textasciicircum uþær ÷ þæ¶r ÷ l-{}-{}-r(a) ÷ (þ)en-nþþæþeskulhuærki} \\
§ D \textbf{skulæ ÷ huærki ÷ søtæn ÷ sofæ ÷ æþ ÷ uarmnæn ÷ uakæ ÷ førr æn ÷ þu ÷ þæssa ÷ bot ÷ biþær ÷ þær ÷ a\textasciicircum k o\textasciicircum rþ ÷ at kæþæ ÷ ro\textasciicircum nti ÷ amæn ÷ o\textasciicircum k þæt ÷ se +}

\sectionline

\newpage

\subsection{Text}

\bvg\bva[]%
\edtrans{\alst{Jo}rð bið ak varðę \hld\ ok \alst{u}p-himęn}{I bid Earth to guard and Up-heaven}{\Bfootnote{\emph{Jorð \dots\ ok up-himęn} ‘Earth \dots\ and Up-heaven’ is a formulaic and originally pre-Christian merism representing the totality of the cosmos as separated from the primæval ettin-hermaphrodite Yimer; cf.~\textlink{Voluspa}[3].  It is here used in a syncretic magical context; the galer (galder-singer) asks the Cosmos itself to watch over and protect him/her as he/she utters the magic words.
A notable parallel to the present instance is found in the OE \textlink{Acerbot}[1]/4 (ca.~1000 \ce), edited above.  Cf.~also a 14th c.~version of the Icelandic \emph{Griða-mál} ‘Grith-speech’, a legal formula for the establishments of truces between former enemies, where we read that: \emph{Jǫrð rę́ðr griðum fyr neðan, en upp-himinn varðar fyr ofan, en sjór fyr útan, sá er kringir um ǫll lǫnd.} ‘The earth rules the griths from below, and Upheaven guards them from above, and the sea from without which encircles all lands.’ \parenciteshorthand[659]{DI_2}.
Due to its invocation of pre-Christian cosmology along with later use of pre-Christian legal formulae like \emph{rę́kr ok rekinn} and \emph{grid-nídings nafn} it is likely that the \emph{Grida-mál} reflects a late remnant of pre-Christian law, and together with the present galder and \textlink{Acreboot}[1]/4 all three texts appear to reflect a common pre-Christian Germanic ritual expression.}} &
\alst{s}ól ok \alst{s}antę María \hld\ ok \alst{s}alfęn Guð dróttęn &
þęt hann \alst{l}ę́ mik \alst{l}ę́knęs-hand \hld\ ok \alst{l}yf-tungę &
at lyfę \emph{\alst{b}}ifj\emph{a}ndę \hld\ þęr \alst{b}ótę þarf.\eva

\bvb I bid Earth to guard and Up-heaven, \\
the Sun and Saint Mary, and the very Lord God \\
that he lend me a leecher’s hand and a medicine-tongue \\
as medicine for the trembler who needs the cure.\evb\evg


\bvg\bva[][5]%
\ind Ór \alst{b}ak ok ór \alst{b}rýst &
\ind ór \alst{l}íkę ok ór \alst{l}im &
\ind ór \alst{ø̂}vęn ok ór \alst{ø̂}ręn &
ór \alst{a}llę þé \hld\ þęr \alst{i}llt kann í at·kumę.\eva

\bvb Out of back and out of breast! \\
Out of body and out of limb! \\
Out of eyes and out of ears! \\
Out of everything, where evil might come in!\evb\evg


\bvg\bva[][9]%
Svart hêtęr \alst{st}ênn \hld\ hann \alst{st}ę́r í hafę útę, &
\ind þęr liggęr á þé \alst{n}í\emph{u} \alst{n}auðęr; &
\ind \edtext{þę́}{\Afootnote{add., but scraped off and poorly legible \emph{\textbf{†÷ þæ¶r ÷ l-{}-{}-r(a) ÷ (þ)en-nþþæþeskulhuærki†}} (a poorly carved form of ll.~9--10) \emph{DR EM85;493}}} skulę hvęrki \alst{s}ǿtęn \alst{s}ofę; &
\ind ęð \alst{v}armęn \alst{v}akę; &
\ind førr ęn þú þęssa \alst{b}ót \alst{b}iðęr, &
þęr \alst{a}k \alst{o}rð \hld\ \alst{a}t·k\emph{v}ę́ðę r\emph{ú}nti. &
\ind Amęn ok þat sé.\eva

\bvb Swart is a stone called; it stands out in the ocean. \\
\ind There lie on it nine needs; \\
\ind they will neither sleep sweetly \\
\ind nor wake warmly, \\
\ind until thou prayest this cure \\
to which I have spoken the runic(?) words. \\
\ind Amen and let it be so.\evb\evg

\sectionline

\newpage
\input{books/galders/Canterbury galder.tex}
\newpage
\input{books/galders/Sigtuna rib.tex}
\newpage
\input{books/galders/Sigtuna Plate I.tex}
\newpage
\input{books/galders/Bryggen.tex}


%\section{Runic plates}
